\documentclass[12pt,a4paper]{article}
\usepackage{cmap}
\usepackage[utf8]{inputenc}
\usepackage[T1,T2A]{fontenc}
\usepackage[english,russian]{babel}
\usepackage{amsmath,booktabs,array}
\usepackage{caption,float,qrcode}
\usepackage{geometry,hyperref,indentfirst}
\geometry{margin=25mm}
\hypersetup{unicode=true,colorlinks=true,urlcolor=blue,linkcolor=black}
\setlength{\parindent}{1.25cm}
\setlength{\parskip}{0pt}
\setlength{\emergencystretch}{2em}
\DeclareCaptionLabelSeparator{emdash}{\space---\space}
\captionsetup[figure]{justification=centering,singlelinecheck=false,labelsep=emdash}
\addto\captionsrussian{\renewcommand{\figurename}{Рис.}}
\linespread{0.96}
\makeatletter
\newcommand{\inputrows}[1]{\@@input #1 }
\makeatother
\newcommand{\LabElevenRawMB}{4140,27}
\newcommand{\LabElevenCompactMB}{1210,54}
\newcommand{\LabElevenSavings}{70,8}
\newcommand{\LabElevenDocs}{478\,833}
\newcommand{\LabElevenTerms}{2\,545\,616}
\newcommand{\LabElevenVerifySec}{86,0}
\newcommand{\LabElevenWholeRawMB}{4443,34}
\newcommand{\LabElevenPackedMB}{1400,56}
\newcommand{\LabElevenWholeSavings}{68,5}
\newcommand{\LabElevenVocabRawMB}{138,83}
\newcommand{\LabElevenVocabCompactMB}{25,77}
\newcommand{\LabElevenVocabSavings}{81,4}

\newcommand{\ProfileRepeats}{7}
\newcommand{\ProfileGodRaw}{127,19}
\newcommand{\ProfileGodVB}{166,87}
\newcommand{\ProfileGodSlower}{31,2}

\begin{document}
\hypersetup{pageanchor=false}
\begin{titlepage}
\begin{center}
{\large МОСКОВСКИЙ АВИАЦИОННЫЙ ИНСТИТУТ\\
(НАЦИОНАЛЬНЫЙ ИССЛЕДОВАТЕЛЬСКИЙ УНИВЕРСИТЕТ)}
\vspace{18pt}
{\large Факультет информационных технологий и прикладной математики}

{\large Кафедра вычислительной математики и~программирования}
\vfill
{\large\textbf{Лабораторная работа №11 по курсу\\ «Информационный поиск»}}
\vspace{12mm}

{\large\textbf{Сжатие индекса}}
\vfill
\begin{flushright}
\begin{tabular}{@{}r@{:\hspace{0.6em}}l@{}}
\textit{Студент} & К.\,А. Арусланов\\
\textit{Группа} & М8О-403Б-23\\
\textit{Преподаватель} & А.\,А. Кухтичев\\
\textit{Дата} & 08.11.2026\\
\textit{Оценка} &
\end{tabular}
\end{flushright}
\vfill
Москва, 2026
\end{center}
\end{titlepage}
\hypersetup{pageanchor=true}
\setcounter{page}{2}
\tableofcontents
\newpage

\section{Постановка задачи}
Разработать способ сжатия инвертированного индекса без потери информации.
Обеспечить чтение сжатых списков при выполнении поисковых запросов и проверить
совпадение выдачи с исходным индексом. Измерить размер данных, время поиска,
чтения и декодирования списков. Подготовить консольную демонстрацию и автотесты.
Основные алгоритмы и структуры данных реализовать на C++.

\section{Метод решения}
\subsection{Состав индекса и область сжатия}
В инвертированном индексе различаются словарь терминов, списки документов
(постинги), координаты вхождений и сведения о самих документах.
Исходный корпус текстов и таблица документов относятся к другим данным.
В данной работе сжаты словарь терминов, постинги и координаты.
Таблица документов не менялась. Морфологический словарь лемматизатора хранится
отдельно и не входит в измеряемый индекс. Экономия оценивается для каждой
структуры и для всего рабочего индекса.

Булев индекс хранит для каждого термина возрастающие идентификаторы документов
в виде 64-битных целых. Координатное расширение хранит пары из порядкового номера
документа и номера токена с признаком поля; каждая пара занимает 8 байт.
Списки упорядочены, поэтому соседние значения часто различаются на небольшую величину.

\subsection{Кодирование разностей и чисел}
Для списка $d_1<d_2<\cdots<d_n$ записывается $\Delta_i=d_i-d_{i-1}$,
где $d_0=0$. Например, список $(3,8,42)$ превращается в $(3,5,34)$.
Для координаты $(o,f,p)$ отдельно кодируются $o-o_{\mathrm{prev}}$ и
$2\Delta p+f$, где $f\in\{0,1\}$ обозначает заголовок или основной текст.
При смене документа либо поля $\Delta p=p$, иначе $\Delta p=p-p_{\mathrm{prev}}$.
Так сохраняется строгий порядок координат и не смешиваются позиции разных полей.

Неотрицательные числа записываются кодом Variable Byte (VB) в варианте лекции.
В каждом байте семь бит данных; младшие группы идут первыми. Старший бит равен
нулю у промежуточных байтов и единице у последнего. Поэтому 128 записывается
двумя байтами \texttt{00 81}. Значения до 127 занимают один байт, до
16\,383 --- два. Декодер
отклоняет усечённые, избыточные и переполненные представления.

\subsection{Блочное сжатие словаря}
Термины упорядочены лексикографически внутри двух пространств. Они разбиты на
блоки по 16 записей; для каждого блока хранится одно смещение. Первый термин
записывается полностью, каждый следующий --- длиной общего префикса с предыдущим
и собственным суффиксом. Например, для \texttt{game} и \texttt{games} вторая
строка представляется длиной префикса 4 и суффиксом \texttt{s}.
В блоке также закодированы пространство термина, число документов и начальное
смещение постингов. Поиск бинарно находит блок по первому термину и последовательно
просматривает до 16 записей. Полный словарь в память не загружается.

\subsection{Формат и чтение}
Для каждого термина сохраняются смещение и длина его сжатого списка.
Это позволяет открыть нужный поток без загрузки всех списков в память.
Сжатое расширение содержит четыре файла числовых потоков и две структуры
словаря: таблицу смещений блоков и поток префиксных записей.
Заголовки имеют номер формата, размеры,
контрольные суммы FNV-64 и отпечаток исходного индекса. Незавершённая сборка
помечается отдельно. Сжатие не меняет токенизацию, лемматизацию, язык запросов
или порядок выдачи.

Для проверки физической экономии подготовлена отдельная рабочая копия индекса.
В ней сохраняются таблица документов и таблица терминов координатного
расширения. От исходных таблицы терминов и двух больших числовых потоков
остаются 80-байтные заголовки для привязки форматов. Из общего файла строк
удаляются строки терминов, а строки заголовков, URL и источников сохраняются.
Поиск в рабочей копии не требует исходных строк терминов, массивов постингов
и координат. Исходный индекс хранится отдельно как эталон сравнения.

\section{Журнал выполнения}
Сначала построены потоки разностей по проверенным исходным индексам.
Из-за произвольных 64-битных ID документов нельзя кодировать усечённые
32-битные идентификаторы; используются полные 64-битные разности.
Координаты сжимаются отдельно от списков документов, поскольку цитатный поиск
должен восстанавливать порядок и соседство слов без потери информации.
При проверке самостоятельной копии обнаружилось, что маркер незавершённой
сборки находился в её корне, а поиск проверял только вложенные каталоги.
Проверка корневого маркера добавлена в читатель сжатого словаря и закреплена
отдельным тестом: незавершённая копия отклоняется, после снятия маркера
готовая копия снова открывается.

Для разбора замедления фразового поиска добавлен диагностический режим.
Один и тот же VB-декодер работает с файлом и с уже загруженным массивом байтов,
что позволяет отдельно измерить чтение и восстановление значений.

\section{План и результаты тестирования}

План тестирования включает:
\begin{enumerate}
\item проверки граничных значений переменного кода, избыточных форм и усечения;
\item построение на малом корпусе, сравнение каждого термина, ID и координаты;
\item сравнение булевых и фразовых запросов в пространствах словоформ и лемм,
включая повторные слова, отрицание и 64-битные ID;
\item отказ при повреждении, неполной сборке, чужом индексе, повторной записи
в существующий каталог и недостатке свободного места;
\item проверку полного корпуса и повторные измерения запросов.
\end{enumerate}
Нативные проверки реализованы в GoogleTest; запуск и данные измерений приложены
к исходникам. Программа \texttt{compact\_index verify-packed} сверяет контрольные
суммы, границы списков и \emph{каждый} восстановленный термин, ID и координату
с исходными файлами.
Пройдены все девять проверок группы ЛР11 и 123 проверки общего набора CTest.

\section{Результаты на корпусе}
На корпусе из \LabElevenDocs{} документов построен индекс с
\LabElevenTerms{} терминами в двух пространствах. Исходные списки документов
и координаты занимают \LabElevenRawMB{} МБ. Четыре сжатых файла вместе
с таблицами смещений занимают \LabElevenCompactMB{} МБ: экономия
\LabElevenSavings\,\%. Полное сравнение декодирования выполнено за
\LabElevenVerifySec{} с.
Полный исходный индекс, включая словарь и метаданные документов, занимает
\LabElevenWholeRawMB{} МБ. Отдельная рабочая копия со сжатыми структурами занимает
\LabElevenPackedMB{} МБ; фактическая экономия для всего индекса составляет
\LabElevenWholeSavings\,\%.

\begin{center}
\begin{tabular}{@{}lrr@{}}
\toprule
Структура & Исходный размер, МБ & Сжатый размер, МБ\\
\midrule
Списки документов & 1\,455{,}88 & 224{,}28\\
Координаты & 2\,684{,}39 & 904{,}80\\
Таблицы смещений & --- & 81{,}46\\
Словарь терминов & \LabElevenVocabRawMB{} & \LabElevenVocabCompactMB{}\\
\bottomrule
\end{tabular}
\end{center}
Словарь терминов уменьшился на \LabElevenVocabSavings\,\% благодаря хранению
общих префиксов, переменному кодированию метаданных и одному смещению
на блок из 16 терминов. Здесь и далее МБ означает $10^6$ байт.
Исходные постинги хранят произвольные 64-битные ID документов, поэтому
процент экономии относится именно к этому формату. Прямое сравнение
с лекционными замерами для 32-битных ID было бы некорректным.

Для набора булевых и фразовых запросов в обоих пространствах совпали полное число
найденных документов и показанные ID. Полная проверка списков отдельно
подтвердила точное восстановление данных. Порядок выдачи остался
по возрастанию ID.

\subsection{Время поиска}
Для дополнительного эксперимента выбраны фразы и булев запрос из тех же слов.
В таблице приведены медианы \ProfileRepeats{} повторений после одного
прогревочного пакета запросов. Сжатый режим использует самостоятельную
рабочую копию. Время включает разбор запроса и перечисление всех совпадений,
но исключает открытие индекса, загрузку морфологического словаря и вывод
метаданных. Кэш ОС не очищался, поэтому это не измерение холодного диска.

\begin{center}
\small
\begin{tabular}{@{}llrr@{}}
\toprule
Пространство & Запрос & Исходный, мс & VB, мс\\
\midrule
\inputrows{lab11/profile/query_rows.tex}
\bottomrule
\end{tabular}
\end{center}
Запросы в кавычках выполнялись как фразы, а строки с AND --- как
пересечение списков документов без проверки соседства слов.
Фраза «god of war» по словоформам выполнялась за \ProfileGodRaw{} мс
на исходном индексе и за \ProfileGodVB{} мс на сжатом.
Сжатие уменьшило размер данных, но в этом запросе время выросло
на \ProfileGodSlower\,\%.

\subsection{Чтение и декодирование}
Отдельно измерены полные списки терминов \texttt{god}, \texttt{of} и
\texttt{war}. Чтение --- это \texttt{seek} и \texttt{fread} в заранее
выделенный и заполненный буфер. Декодирование выполняется из памяти без
обращений к файлу: для исходного формата читаются фиксированные целые,
для VB восстанавливаются разности и поля координат. Оба варианта вычисляют
одинаковую контрольную сумму восстановленных значений.
Результаты сравниваются после замера; расхождений не обнаружено.
Порядок исходного и сжатого вариантов чередуется между повторениями.

В таблице показаны координаты в пространстве словоформ, время дано в мс.
Для каждого списка выполнен один предварительный прогрев; приведены
медианы \ProfileRepeats{} повторений.
\begin{center}
\begin{tabular}{@{}lrrrr@{}}
\toprule
Термин & \multicolumn{2}{c}{Чтение} & \multicolumn{2}{c}{Декодирование}\\
 & Исходный & VB & Исходный & VB\\
\midrule
\inputrows{lab11/profile/decode_rows.tex}
\bottomrule
\end{tabular}
\end{center}

Для связи с реальным поиском также измерен полный проход штатного курсора:
открытие потока, чтение и восстановление всех координат, вычисление той же
контрольной суммы и закрытие потока.
\begin{center}
\begin{tabular}{@{}lrr@{}}
\toprule
Термин & Исходный курсор, мс & VB-курсор, мс\\
\midrule
\inputrows{lab11/profile/cursor_rows.tex}
\bottomrule
\end{tabular}
\end{center}

Особенно заметна разница на \texttt{of}: список содержит 1\,997\,146 координат
и занимает 15{,}98 МБ в исходном формате против 4{,}37 МБ в VB.
Сжатые байты читаются быстрее, но их восстановление требует больше времени.
Штатный курсор также выполняет побайтовое чтение VB и проверки каждой
координаты. Это согласуется с замедлением фразового запроса на прогретых данных.
Для списков ID картина другая: уменьшение числа читаемых байтов может
перевесить затраты на более простое восстановление разностей.

Эти полные проходы являются отдельной нагрузкой. Фразовый поиск может читать
лишь часть списков и дополнительно сопоставляет координаты, поэтому времена
проходов нельзя складывать для получения времени запроса.
Исходные ответы, все повторения и замеры обоих пространств сохранены
в \path{lab11/profile/}; запуск описан в \path{lab11_run.md}.

Исходный код:
\href{https://github.com/kirbasss/MAI-IR/tree/lab11}{github.com/kirbasss/MAI-IR/tree/lab11}.
\begin{figure}[H]
\centering
\begingroup
\hypersetup{urlcolor=black}
\qrcode[height=25mm]{https://github.com/kirbasss/MAI-IR/tree/lab11}
\endgroup
\caption{Ссылка на ветку лабораторной работы №11}
\end{figure}

\newpage
\section{Вывод}
В этой работе я добавил сжатие словаря, списков документов и координат слов.
Полный рабочий индекс удалось уменьшить с \LabElevenWholeRawMB{} до
\LabElevenPackedMB{} МБ, то есть на \LabElevenWholeSavings\,\%.
Проверка всех восстановленных данных не выявила отличий от исходного индекса,
а на тестовых запросах результаты поиска совпали.

Лучше всего сократились списки документов: вместо полных ID в них хранятся
небольшие разности. При этом сжатие не всегда ускоряет поиск.
Например, запрос «god of war» по словоформам стал выполняться дольше.
Отдельные замеры показали, что сжатые данные читаются быстрее,
но восстановление координат требует больше времени.
Получился заметный выигрыш по месту, но скорость ещё есть куда улучшать.

\end{document}
